\chapter{Introdução}
\label{cap:introducao}

\section{Contextualização e motivação}
\label{sec:contextualizacao}

% Tópicos sugeridos:
%   - por que moléculas fotoativas/funcionais importam na prática
%     (ex.: OLEDs, células solares, singlet fission, corantes/fluoróforos)
%   - lacuna entre desenho racional de moléculas e exploração computacional
%     do espaço químico
%   - papel do SMILES como representação computacional de moléculas

\section{Definição do problema}
\label{sec:problema}

% Tópicos sugeridos:
%   - o que exatamente se quer gerar/otimizar (classe de moléculas-alvo)
%   - por que geração combinatória de fragmentos é o ponto de partida
%   - como diversidade/redundância estrutural entra como problema em aberto

\section{Organização do texto}
\label{sec:organizacao}

% TODO: um parágrafo curto remetendo a cada capítulo, escrito por último.
